\subsection{Inverse of a Relation} 
Every relation has an \term{inverse}
\begin{definition}[Inverse of a Relation]
For a relation $R$, its inverse, written $R\inv$, is the relation
\begin{center}$R\inv=\Set{(x,y)\mid (y,x)\in R}$\end{center}
Read $R\inv$ as ``R inverse.''
\end{definition}
\begin{note}The ${}\inv$ in $R\inv$ should be read as ``inverse.'' It looks like an exponent, but it isn't one.\end{note}
\begin{example} Give the inverse of each relation.
\begin{enumerate}\setabc
\item $R=\Set{(-9,1),(-5,1),(-2,9),(-2,8)}$\lbr$R\inv=\Set{\makeblank[3in]}$
\item $S=\Set{(-2,-3),(4,4),(6,8),(6,-5)}$\lbr$S\inv=\Set{\makeblank[3in]}$
%\item $T=\Set{(-3,-9),(5,-1),(6,7),(9,4)}$\lbr$T\inv=\Set{\makeblank[3in]}$
%\item $Q=\Set{(-6,-3),(-1,9),(4,2),(7,-3)}$\lbr$Q\inv=\Set{\makeblank[3in]}$
\end{enumerate}
\end{example}
Reversing the pairs also reverses the \makeblank[1in] and \makeblank[1in].
\begin{example}
Suppose $R$ is a relation with $\dom R=\Set{1,3,5}$ and $\rng R=\Set{-2,2}$.
\begin{itemize}
\item $\dom R\inv=\makeblank[1in]=\Set{\makeblank[1in]}$
\item $\rng R\inv=\makeblank[1in]=\Set{\makeblank[1in]}$
\end{itemize}
\end{example}
\begin{example}
Let $f(x)=x^2,\ x\geq 0$. What is $\rng f\inv$?\makeblanknewf\makeblanknewf
\end{example}
\begin{example}
Suppose $h$ is a one-to-one function, with $\dom h=\Set{x\mid x\neq 4}$ and\linebreak$\dom h\inv=\Set{x\mid x\neq -2}$.
\begin{itemize}
\item $\rng h=\makeblank[1in]=\Set{\makeblank[1in]}$
\item $\rng h\inv=\makeblank[1in]=\Set{\makeblank[1in]}$
\end{itemize}
\end{example}